\cleardoublepage
\thispagestyle{empty}

\begin{center}
\textbf{DECLARAÇÃO DE USO DE INTELIGÊNCIA ARTIFICIAL GENERATIVA}
\end{center}

\vspace{1.5cm}

Durante a elaboração deste Trabalho de Conclusão de Curso, foram utilizadas ferramentas de Inteligência Artificial Generativa (IAG). Esta declaração identifica as ferramentas e as finalidades efetivas do apoio, em consonância com a orientação de transparência da Portaria CNPq nº 2.664/2026 \cite{cnpq2026}.

\textbf{Ferramentas utilizadas:} Claude (Anthropic) e Codex (OpenAI), modelos de linguagem de propósito geral.

\textbf{Finalidade do uso:} as ferramentas apoiaram a reorganização do tema, a estruturação e revisão textual, o levantamento e a conferência de fontes, a formatação, a elaboração de diagramas, a discussão do protocolo, a geração e revisão de código, a preparação e execução assistida de testes e da campanha, a agregação dos dados e a redação dos resultados e conclusões. O apoio incluiu comandos e scripts de análise; não ficou restrito à correção gramatical.

As medições apresentadas foram obtidas pela execução do protótipo e dos instrumentos, e não por simulação textual das ferramentas de IAG. Os dados brutos e os vereditos originais foram preservados, inclusive a evidência temporal incompleta. A análise exploratória posterior foi identificada separadamente e não substituiu os critérios originais.

O autor assume responsabilidade pelo conteúdo, pelas fontes, pelo método, pelo código, pelos dados e pelas conclusões, inclusive pelos trechos elaborados com auxílio de IAG. Na revisão de 30 de outubro de 2026, informou que a conferência pessoal integral do código, das decisões metodológicas, das tabelas e das conclusões ainda seria realizada antes da entrega. Esta versão não declara essa conferência como concluída nem presume aprovação da orientadora.

\cleardoublepage
